% Chapter 1

\chapter{Introducción general} % Main chapter title

\label{Chapter1} % For referencing the chapter elsewhere, use \ref{Chapter1} 
\label{IntroGeneral}

En este capítulo se describen las principales características de la demanda eléctrica en la República Argentina y los desafíos que presenta pronosticar su comportamiento. Además, se pone de manifiesto su relevancia en la planificación y operación del Sistema Argentino de Interconexión (SADI).

%----------------------------------------------------------------------------------------

% Define some commands to keep the formatting separated from the content 
\newcommand{\keyword}[1]{\textbf{#1}}
\newcommand{\tabhead}[1]{\textbf{#1}}
\newcommand{\code}[1]{\texttt{#1}}
\newcommand{\file}[1]{\texttt{\bfseries#1}}
\newcommand{\option}[1]{\texttt{\itshape#1}}
\newcommand{\grados}{$^{\circ}$}

%----------------------------------------------------------------------------------------

%\section{Introducción}

%----------------------------------------------------------------------------------------


\section{Introducción}
\label{sec:introduccion}

La energía eléctrica constituye un insumo esencial para el funcionamiento de los hogares, la industria, el comercio, los servicios y la infraestructura de un país. Su disponibilidad debe garantizarse de manera continua, con niveles adecuados de calidad y seguridad. Dado que no puede almacenarse en grandes cantidades, la generación debe ajustarse permanentemente al consumo, junto con las pérdidas y los intercambios que se producen en la red.

La demanda eléctrica representa el consumo agregado de los usuarios conectados al sistema. Su comportamiento varía a lo largo del día según la actividad residencial, comercial e industrial. También presenta diferencias entre días hábiles, fines de semana y feriados, así como variaciones estacionales asociadas, entre otros factores, con las condiciones meteorológicas. Por lo tanto, constituye una variable temporal cuyo comportamiento depende de múltiples factores.

En la Argentina, el abastecimiento eléctrico se realiza principalmente a través del Sistema Argentino de Interconexión (SADI), que vincula centrales de generación, redes de transporte, instalaciones de distribución y grandes usuarios. Su operación requiere coordinar recursos con distintas características técnicas y económicas, entre ellos centrales térmicas, hidroeléctricas, nucleares y fuentes renovables variables. También deben considerarse los intercambios con países limítrofes, las restricciones de la red, las reservas operativas y la disponibilidad de combustibles y recursos hidráulicos.

La planificación de la operación permite anticipar estas condiciones y definir una estrategia para abastecer la demanda esperada. La programación semanal contempla decisiones relacionadas con la disponibilidad de generación, los mantenimientos, el uso de los recursos hidráulicos, la asignación de combustibles y los intercambios de energía. La programación diaria establece una previsión más detallada del funcionamiento del sistema para el día siguiente.

En ambas instancias, una subestimación de la demanda puede conducir a una programación insuficiente de recursos y reducir los márgenes disponibles para responder ante contingencias. Por el contrario, una sobreestimación puede provocar la programación innecesaria de unidades de generación e incrementar los costos operativos. En consecuencia, la calidad del pronóstico influye en la confiabilidad, la seguridad y la eficiencia económica de la operación.

El pronóstico de la demanda consiste en estimar su evolución futura a partir de información histórica y de variables relacionadas con su comportamiento. Los pronósticos de corto plazo se utilizan principalmente en la programación y operación del sistema, mientras que los horizontes más extensos permiten evaluar inversiones, ampliaciones de infraestructura y escenarios de crecimiento. La naturaleza temporal de la demanda, su sensibilidad a factores meteorológicos y calendarios, y la necesidad de anticipar su evolución justifican el desarrollo de métodos capaces de representar relaciones complejas y adaptarse a diferentes condiciones operativas.

%----------------------------------------------------------------------------------------

%\section{Contexto y motivación}

%----------------------------------------------------------------------------------------


\section{Contexto y motivación}
\label{sec:contexto_motivacion}

El trabajo se realizó para la Compañía Administradora del Mercado Mayorista Eléctrico S.A. (CAMMESA), en particular para el área de Programación Semanal y Diaria de la Gerencia de Coordinación Operativa. Esta área requiere estimaciones de demanda con un horizonte de hasta catorce días para la programación semanal y un perfil horario para la programación del día siguiente.

Antes de la realización del trabajo, el área disponía de modelos de machine learning y deep learning destinados a pronosticar la demanda diaria del SADI. Sin embargo, estos modelos no contaban con un proceso sistemático de validación y monitoreo que permitiera evaluar su desempeño a lo largo del tiempo y bajo diferentes condiciones meteorológicas y calendarias. En consecuencia, no se disponía de información suficiente para determinar objetivamente la confiabilidad de las predicciones ni para detectar una posible degradación de su rendimiento.

Tanto en la programación semanal como en la diaria, los resultados de los modelos se utilizaban como referencia y se complementaban mediante la inspección manual de datos históricos de demanda y temperatura. Los especialistas analizaban períodos con características similares e inferían los valores esperados para los días siguientes. Este procedimiento requería revisar distintas fuentes de información y dependían en gran medida de la experiencia de los especialistas.

La ejecución de los modelos también presentaba dificultades operativas, ya que se encontraban funcionando en una única computadora de uso común y el flujo de datos tenía un bajo grado de automatización. Su utilización requería cargar datos de entrada, ejecutar macros para obtener o transformar información e iniciar manualmente los procesos correspondientes. Tampoco existía un mecanismo centralizado para almacenar las ejecuciones, las variables de entrada y las predicciones generadas, lo que dificultaba la comparación entre modelos, el cálculo periódico de métricas y el análisis de los errores.

Otra limitación correspondía al nivel de agregación geográfica. Los modelos diarios generaban pronósticos para la totalidad del SADI, pero no proporcionaban estimaciones por región eléctrica. CAMMESA también dispone de un modelo de demanda horaria limitado al Gran Buenos Aires, cuyos resultados pierden confiabilidad al extenderse el horizonte de predicción. Por lo tanto, las herramientas existentes no permitían obtener pronósticos diarios regionales para el horizonte requerido por la programación semanal.

En este contexto, el trabajo se orientó a obtener una herramienta de pronóstico confiable y adecuada para su utilización operativa. Se automatizaron la obtención y el procesamiento de los datos, la ejecución de las predicciones y el almacenamiento de los resultados. También se incorporaron mecanismos de validación y monitoreo basados en métricas objetivas, se compararon diferentes modelos y se analizó la generación de pronósticos por región eléctrica.

Esta memoria describe el desarrollo y la evaluación de la solución implementada, que permite reducir la intervención manual, mejorar la trazabilidad y disponer de información objetiva sobre la calidad de las predicciones. De este modo, los especialistas pueden concentrarse en el análisis de situaciones particulares y utilizar los resultados como apoyo para la programación semanal y diaria del sistema eléctrico.


%----------------------------------------------------------------------------------------

%\section{Estado del arte}

%----------------------------------------------------------------------------------------

\section{Estado del arte}
\label{sec:estado_arte}

El pronóstico de la demanda eléctrica ha sido abordado mediante modelos estadísticos, técnicas de aprendizaje automático y arquitecturas de aprendizaje profundo. La elección del método depende del horizonte de predicción, la resolución temporal, el nivel de agregación geográfica y la disponibilidad de variables explicativas. En sistemas con una participación residencial significativa, la temperatura y otras condiciones meteorológicas tienen una incidencia relevante sobre el consumo, principalmente por el uso de equipos de climatización. En un país de gran extensión territorial, este efecto varía entre regiones y estaciones del año, por lo que un pronóstico agregado puede ocultar comportamientos locales.

En la Argentina, Oviedo desarrolló un modelo LSTM para pronosticar la demanda horaria de la provincia de Córdoba a partir de datos de consumo y variables climáticas, temporales y demográficas \cite{oviedo2025regional}. El modelo alcanzó un error porcentual absoluto medio (MAPE) de 3,20\% y un coeficiente de determinación de 0,95. Por su parte, Uhrig et al. estudiaron la demanda diaria de Entre Ríos mediante una LSTM multivariada con variables meteorológicas, temporales y energéticas \cite{uhrig2026multivariate}. Los autores observaron mayores niveles de demanda ante temperaturas extremas y obtuvieron mejores resultados que con un método estadístico y una red GRU. Ambos antecedentes confirman la relevancia de las variables exógenas, aunque se limitan a una provincia y no consideran conjuntamente la demanda nacional y su desagregación regional.

Como antecedente latinoamericano de alcance nacional, Leme et al. aplicaron bosques aleatorios, máquinas de vectores de soporte y potenciación por gradiente al pronóstico de la demanda del sistema interconectado brasileño \cite{leme2020brazil}. Aquila et al. revisaron los métodos utilizados para la planificación operativa de corto plazo en Brasil y destacaron la selección de variables, la segmentación según patrones similares y los enfoques jerárquicos de pronóstico \cite{aquila2023overview}. Debido a su extensión territorial y diversidad climática, el caso brasileño presenta mayor proximidad con el problema argentino que los estudios limitados a una empresa distribuidora o a grupos reducidos de usuarios.

La experiencia operativa de CAMMESA constituye otro antecedente directo del trabajo. Antes de su realización, el área de Programación Semanal y Diaria utilizaba modelos SARIMAX, XGBoost y LSTM para estimar la demanda agregada del SADI. Las evaluaciones internas mostraban valores de MAPE aproximadamente comprendidos entre 3\% para el día siguiente y 10\% para los horizontes más alejados, con diferencias según el modelo y las condiciones analizadas. Estos resultados no proceden de una publicación académica y, por lo tanto, no son directamente comparables con los trabajos citados. No obstante, documentan el desempeño de los métodos en el entorno de aplicación y muestran tanto el aumento del error con el horizonte como la necesidad de incorporar variables meteorológicas y calendarias identificadas por los especialistas.

Los métodos basados en árboles, como XGBoost, representan una referencia adecuada para este problema porque permiten modelar relaciones no lineales e incorporar variables exógenas con un costo computacional moderado. Sin embargo, requieren construir explícitamente atributos que describan la historia y la estacionalidad de la serie. En el trabajo, XGBoost se utilizó como línea de base para determinar si enfoques más recientes mejoraban el pronóstico bajo un mismo procedimiento de validación.

La investigación reciente avanzó hacia arquitecturas neuronales diseñadas para generar horizontes de múltiples pasos y representar patrones en distintas escalas. Oreshkin et al. aplicaron N-BEATS al pronóstico de demanda eléctrica de mediano plazo sobre 35 series nacionales europeas y obtuvieron mejoras frente a métodos estadísticos, de aprendizaje automático e híbridos \cite{oreshkin2021nbeatsload}. El resultado aporta evidencia específica sobre la utilidad de las arquitecturas residuales para modelar tendencias, estacionalidades y fluctuaciones de la demanda, aunque trabaja con series mensuales y un horizonte diferente del considerado en este trabajo. Otras arquitecturas, entre ellas N-HiTS y modelos basados en Transformers como PatchTST, extienden esta línea mediante mecanismos de procesamiento multiescala y segmentación temporal \cite{challu2023nhits,nie2023patchtst}. Su interés para el problema radica en la capacidad de producir pronósticos directos de varios pasos y aprovechar ventanas históricas extensas, pero su desempeño debe comprobarse con datos meteorológicos y regionales del sistema analizado.

Los modelos fundacionales constituyen una tendencia aún más reciente. Se preentrenan sobre grandes colecciones de series y pueden aplicarse a datos nuevos sin entrenamiento específico o con un ajuste limitado. Meyer et al. compararon modelos fundacionales, entre ellos Chronos y TimesFM, con Transformers entrenados desde cero para pronosticar demanda eléctrica residencial de corto plazo \cite{meyer2025foundationload}. Los resultados indicaron desempeños comparables y, en algunos casos, superiores, con menor necesidad de entrenamiento específico. Este antecedente respalda su evaluación en el dominio eléctrico, aunque su escala residencial difiere de la demanda de un sistema nacional y no resuelve por sí mismo la incorporación de variables exógenas ni la desagregación regional.

La tabla \ref{tab:estado_arte} resume los antecedentes de mayor relación con el trabajo. Las métricas no son directamente comparables debido a las diferencias entre conjuntos de datos, horizontes, resoluciones y niveles de agregación.

\begin{table}[htbp]
	\centering
	\caption{Antecedentes relacionados con el pronóstico de demanda eléctrica.}
	\label{tab:estado_arte}
	\footnotesize
	\setlength{\tabcolsep}{3pt}
	\renewcommand{\arraystretch}{1.10}
	
	\begin{tabularx}{\textwidth}{@{}XXXX@{}}
		\toprule
		\textbf{Referencia} &
		\textbf{Enfoque} &
		\textbf{Datos y alcance} &
		\textbf{Resultado o aporte} \\
		\midrule
		
		Oviedo \cite{oviedo2025regional} &
		LSTM &
		Demanda horaria de Córdoba y variables exógenas &
		MAPE de 3,20\% y $R^2$ de 0,95 \\
		
		Uhrig et al. \cite{uhrig2026multivariate} &
		LSTM multivariada &
		Demanda diaria de Entre Ríos y 21 variables &
		Mejora frente a un método estadístico y una GRU \\
		
		CAMMESA, evaluación interna &
		SARIMAX, XGBoost y LSTM &
		Demanda agregada del SADI, hasta catorce días &
		MAPE aproximado entre 3\% y 10\% según horizonte y modelo \\
		
		Leme et al. \cite{leme2020brazil} &
		RF, SVM y potenciación por gradiente &
		Sistema interconectado brasileño &
		Comparación sobre datos de alcance nacional \\
		
		Oreshkin et al. \cite{oreshkin2021nbeatsload} &
		N-BEATS &
		Demanda mensual de 35 países europeos &
		Mejora de precisión y sesgo frente a diez métodos \\
		
		Meyer et al. \cite{meyer2025foundationload} &
		Modelos fundacionales y Transformers &
		Demanda residencial de corto plazo &
		Desempeño competitivo con menor entrenamiento específico \\
		\bottomrule
	\end{tabularx}
\end{table}

La revisión mostró una disponibilidad limitada de estudios que reprodujeran las condiciones del problema considerado. Los antecedentes argentinos publicados se concentran en una provincia, mientras que los trabajos sobre arquitecturas modernas utilizan otras escalas, horizontes o niveles de agregación. Tampoco se encontraron comparaciones que combinaran un horizonte diario de catorce días, variables meteorológicas, demanda total de un sistema eléctrico nacional y desagregación en regiones extensas.


\section{Objetivos y alcance}
\label{sec:objetivos}

El objetivo general del trabajo fue desarrollar e implementar una solución automatizada para pronosticar la demanda eléctrica diaria del Sistema Argentino de Interconexión (SADI), con un horizonte de hasta catorce días. Las variables objetivo fueron la energía diaria y las potencias mínima y máxima, tanto para la totalidad del sistema como para cada una de las regiones eléctricas consideradas.

Para alcanzar este propósito, se definieron los siguientes objetivos específicos:

\begin{itemize}
	
	\item Integrar y procesar los datos históricos de demanda, las variables meteorológicas y la información calendaria necesarias para generar los pronósticos del SADI y de sus regiones eléctricas.
	
	\item Implementar y comparar modelos de aprendizaje automático, arquitecturas neuronales modernas y modelos preentrenados para series temporales, utilizando un modelo XGBoost como referencia.
	
	\item Evaluar la capacidad de los modelos para pronosticar la energía diaria y las potencias mínima y máxima en cada uno de los catorce días del horizonte.
	
	\item Comparar el desempeño de los modelos mediante el error absoluto medio (MAE), la raíz del error cuadrático medio (RMSE) y el error porcentual absoluto medio (MAPE), calculados sobre períodos no utilizados durante el entrenamiento.
	
	\item Analizar los errores según la variable pronosticada, el horizonte de predicción, la región eléctrica y las condiciones meteorológicas y calendarias.
	
	\item Automatizar la obtención y el procesamiento de los datos, la ejecución periódica de los modelos y el almacenamiento centralizado de sus entradas, predicciones y métricas.
	
	\item Implementar mecanismos de monitoreo que permitan consultar el desempeño histórico de los modelos, comparar sus resultados y detectar variaciones en la calidad de las predicciones.
	
\end{itemize}

El cumplimiento de estos objetivos permitió determinar qué enfoques resultaron más adecuados para las distintas variables, regiones y horizontes, y disponer de una herramienta trazable para apoyar las tareas de programación semanal y diaria de CAMMESA.